% !PRINTED 354
\weberpage{354}
%<<CONT
points other than the points of contact of the double tangents
is satisfied.

The variables $y$ satisfy the tangent equation

\weq{7}{y_1 f'(x_1) + y_2 f'(x_2) + y_3 f'(x_3) = 0.}

If we denote by $b_1, b_2, b_3$ any three arbitrary constants, and set

\weq{8}{b_x = b_1 x_1 + b_2 x_2 + b_3 x_3,}

so that $b_x = 0$ is the equation of an arbitrary straight line, then we may take, in addition to (7), the equation

\weq{9}{b_1 y_1 + b_2 y_2 + b_3 y_3 = 0}

as well, and thus regard $(y)$ as the point of intersection of the tangent at $(x)$ with the arbitrary straight line $b_x$. We then obtain

\weq{10}{\begin{aligned}
y_1 &= b_2 f'(x_3) - b_3 f'(x_2),\\
y_2 &= b_3 f'(x_1) - b_1 f'(x_3),\\
y_3 &= b_1 f'(x_2) - b_2 f'(x_1),
\end{aligned}}

and hereby the condition (7) is identically satisfied. By this substitution $R(x,y)$ passes into a homogeneous function of degree $20$ in the $x$ and of degree $6$ in the $b$, which we shall denote by $D(x,b)$. But this function $D(x,b)$ vanishes not only at the points of contact of the double tangents, but also at the four points of intersection of the line $b_x$ with the curve $f$. This condition must now be transformed further, so that it becomes independent of the $b$.

If we pass from the point $(y)$ to any other point of the tangent, we can achieve this by replacing $y$ by $\mu x + \lambda y$, where $\lambda, \mu$ denote two arbitrary parameters. Thereby $f(x+ty)$ passes into

\[
f[(1+\mu t)x + \lambda t y] = (1+\mu t)^4 f\left(x + \frac{\lambda t}{1+\mu t}\, y\right)
\]

and the development (3) yields, if we set

\[
P_0(x,y),\ P_1(x,y),\ P_1(x,\mu x + \lambda y)
\]

equal to zero:

\[
\begin{aligned}
&6t^2(1+\mu t)^2\lambda^2 P_2(x,y) + 4t^3(1+\mu t)\lambda^3 P_3(x,y) + t^4\lambda^4 P_4(x,y)\\
&= 6t^2 P_2(x,\mu x + \lambda y) + 4t^3 P_3(x,\mu x + \lambda y) + t^4 P_4(x,\mu x + \lambda y),
\end{aligned}
\]
%CONT>>
